% Define el tipo de documento: artículo en este caso.
\documentclass[12pt, a4paper, oneside, titlepage]{article}

%----------- Paquetes necesarios para esta memoria -------------------------------------
\usepackage[utf8]{inputenc}
\usepackage[spanish,es-tabla,es-nodecimaldot]{babel}
\usepackage{amsmath}
\usepackage{amsfonts}
\usepackage{amssymb}
\usepackage{csquotes}	% Manejo de simbolos de cita multiples
\usepackage{graphicx}	% Figuras
% Configuracion de interlineado (\onehalfspace y \singlespacing en capitulos)
%\usepackage{setspace}
\usepackage{float}	% Permite que las figuras floten en la posición en la que aparecen en el código
\usepackage{caption}	% Permite dar estilo a los subtitulos de las figuras (captions)
\usepackage{subcaption}	% Permite sub-captions en sub-figuras (figuras dentro de figuras)
\usepackage{wrapfig}
\usepackage[table]{xcolor}
\usepackage{multirow}	% Permite múltiples filas en tablas
\usepackage{emptypage}	% Elimina numeracion de paginas en blanco
\usepackage{enumitem}	% Permite continuar con la numeración entre listas distintas
\usepackage{algorithm}
\usepackage{algorithmic}

%Paquete para gestionar enlaces
\usepackage{hyperref}
\graphicspath{{./capitulos/img/}} %ruta a la carpeta de imaganes, nos evita poner siempre la ruta

%Paquete para gestionar la bibliografía
\usepackage[backend=bibtex, style=numeric, sorting=none, maxbibnames=6]{biblatex}
%---------------------------------------------------------------------------------------

%Datos del alumno y asignatura
\author{Antonio Ariza García}
\title{Introducción a los modelos computacionales: Práctica 1}

%Fecha de creación de la memoria
\date{Curso académico 2020-2021\\Córdoba, \today}

% Cambia el título del índice de algoritmos
%\renewcommand{\listalgorithmname}{Índice de algoritmos}

% Comienzo del documento
\begin{document}

% Cargar la portada
\begin{titlepage}
\centering
\includegraphics[width=0.18\textwidth]{img/logo_uco_sin_texto.png}
\includegraphics[width=0.30\textwidth]{img/logo_uco_epsc.png}\par\vspace{1cm}
{\scshape\LARGE Universidad de Córdoba\par}
{\scshape\Large Escuela Politécnica Superior de Córdoba\par}
\vspace{1cm}
{\scshape\LARGE Ingeniería Informática\par}
{\scshape\Large Especialidad: Computación\par}
{\scshape\Large Cuarto curso. Primer cuatrimestre\par}
\vspace{1.5cm}
{\scshape\LARGE Introducción a los modelos computacionales.\par}
\vspace{1.5cm}
{\huge\bfseries Práctica 3: Redes neuronales de funciones de base radial.\par}
\vspace{1.2cm}
{\Large\itshape Antonio Ariza García\par}
DNI\par
\href{mailto:correo@uco.es}{i62argaa@uco.es}\par
\vfill

{\large Curso académico 2020-2021\\Córdoba, \today\par}
\end{titlepage}

% Cambia el tipo de numeración de páginas a números romanos
\pagenumbering{roman}
% Muestra el índice de contenido
\tableofcontents
\newpage

% Muestra el índice de figuras y lo incluye en el índice de contenido
\listoffigures
\addcontentsline{toc}{section}{Índice de figuras}
\newpage

% Muestra el índice de tablas y lo incluye en el índice de contenido
\listoftables
\addcontentsline{toc}{section}{Índice de tablas}
\newpage

% Muestra el índice de algoritmos y lo incluye en el índice de contenido
%\listofalgorithms
%\addcontentsline{toc}{section}{Índice de algoritmos}
%\newpage

% Cambia la numeración de páginas a numeración decimal
\pagenumbering{arabic}

% Cambia el espaciado entre párrafos
\setlength{\parskip}{0.7em}

% Cargarmos las demás partes del documento
\section{Introducción} \label{ch:introduccion}

Este documento está realizado en Latex, forma parte de la asignatura Introducción al aprendizaje automático, de la escuela politécnica de Córdoba.
Es una breve introducción a Weka, software utilizado para aprendizaje automático y la minería de datos, está escrito en Java.
\newline
El siguiente documento tiene como objetivo presentarse como guión y entrega de prácticas, puede utilizarse como pequeña guía a la introducción a Weka y aprendizaje automático.

\end{document}